\clearpage
\section{An analytic description of the MC core shift}
The ten accepted native cores at 60--240~MeV support a compact empirical
location law. Four simple forms are fitted with equal weight per mass:
a proportional shift, an affine function, a quadratic and a logarithm.
Each mass is also omitted in turn and predicted from the other nine.
No observed spectrum enters this comparison. The best leave-one-mass-out
(LOO) RMS among these candidates is the logarithmic model,
\begin{equation}
 c(m)-m=-3.22433-2.21399\ln\!\left(\frac{m}{150}\right)\ \mathrm{MeV},
 \qquad 60\leq m\leq240~\mathrm{MeV}.
\end{equation}
The argument uses $m$ in MeV. The law is an empirical approximation,
not a detector mechanism or an extrapolation outside the fitted range.
Its fitted RMS is 0.0497~MeV, LOO RMS 0.0652~MeV and largest held-out
error 0.1604~MeV. These errors describe cross-mass prediction, not a
calibrated confidence band. Model selection and evaluation use the same
ten masses; the LOO summary is a diagnostic, not independent validation.
\fig{.32}{../figures/analytic_center_models.pdf}{Measured core shifts,
analytic models and their held-out residuals. Small error bars show
conditional MC-bin resampling spread; the equal-mass fits do not treat
that spread as the complete response uncertainty.}
\begin{table}[H]\centering\small
\begin{tabular}{lrrrr}
\toprule
Model & Parameters & Fit RMS (MeV) & LOO RMS (MeV) & Max. LOO (MeV) \\
\midrule
proportional & 1 & 0.2552 & 0.2885 & 0.4602 \\
affine & 2 & 0.1555 & 0.2154 & 0.5287 \\
quadratic & 3 & 0.0679 & 0.1399 & 0.3015 \\
logarithmic & 2 & 0.0497 & 0.0652 & 0.1604 \\
\bottomrule
\end{tabular}

\caption{Analytic model comparison. Affine and quadratic forms use
$x=(m-150)/100$; coefficients and all held-out predictions are saved.}
\end{table}

\clearpage
\section{Can one normalized core describe all masses?}
The ten distributions are first recentered without rescaling their
physical widths. They are then expressed in
$u=(m_{ee}-c)/\sigma_{\rm core}$ using the measured local Gaussian width.
The density transformation includes the Jacobian:
$q_m(u)=\sigma_{\rm core}p_m(c+\sigma_{\rm core}u)$.
Each full source histogram has unit area over its recorded 0--400~MeV
range. Display cropping never renormalizes the full distributions.
The separately labeled core view conditions on $|u|<2$ and therefore
discards information about the fraction of all candidates in that core.
\fig{.64}{../figures/shared_shape_overlays.pdf}{All ten normalized shapes
superimposed. Top left: centered physical-mass densities. Top right:
full-normalized densities after width scaling. Bottom left: unit-area
conditional cores. Bottom right: full-normalized tails on a log scale.
The black dashed curve is the equal-mass mean standardized distribution;
sample event counts do not determine its weights.}
The conditional cores are much closer than the full shapes. A shared
central shape is therefore a useful possibility, but its unit-area
overlay cannot establish a common core fraction or common tails. The
full candidate distributions include the same unresolved response and
candidate-association contributions as the earlier sections.

\clearpage
\subsection{Held-out tests of a common shape}
For each generated mass, the common standardized CDF is averaged from
the other nine histograms. The primary shape-only test supplies the held-out
sample's known center and width. A second test predicts its center from
the selected analytic law fitted to the other nine and predicts its
width from a quadratic in $\ln(\sigma_{\rm core}/\mathrm{MeV})$ versus
$x=(m-150)/100$, also fitted without that sample. All likelihood comparisons
use the same measured-center evaluation window and GP constraint as the
direct MC. The second test predicts the complete template but does not
also move its evaluation mask. Existing neighboring-mass morphs are
evaluated in that same window as a comparison.
\begin{table}[H]\centering\scriptsize
\begin{tabular}{rrrrrrr}
\toprule
$m_0$ & Core (\%) & $D_{core}$ & $D_{full}$ & Shared UL ratio & Morph UL ratio & Predictive UL ratio \\
\midrule
60 & 32.4 & 0.045 & 0.461 & 0.443 & 1.047 & 0.419 \\
80 & 60.7 & 0.009 & 0.075 & 0.944 & 1.324 & 0.991 \\
100 & 68.1 & 0.008 & 0.023 & 1.140 & 1.050 & 1.155 \\
120 & 71.0 & 0.010 & 0.061 & 1.191 & 1.025 & 1.172 \\
140 & 71.5 & 0.010 & 0.072 & 1.135 & 1.009 & 1.124 \\
160 & 70.7 & 0.009 & 0.067 & 1.135 & 1.016 & 1.127 \\
180 & 69.6 & 0.007 & 0.052 & 1.085 & 1.017 & 1.076 \\
200 & 67.5 & 0.007 & 0.033 & 1.002 & 1.003 & 1.003 \\
220 & 65.8 & 0.009 & 0.037 & 0.936 & 1.018 & 0.936 \\
240 & 63.6 & 0.014 & 0.049 & 0.864 & 1.005 & 0.860 \\
\bottomrule
\end{tabular}

\caption{Held-out CDF distances and observed upper-limit ratios relative
to direct MC. Core probability is conditional on the stored histogram
range; $D_{core}$ conditions on $|u|<2$. Full distances use the declared
analysis support. ``Predictive'' estimates center and width as well as
shape without the held-out histogram.}
\end{table}
\fig{.26}{../figures/shared_shape_validation.pdf}{Common-shape and
neighbor-morph prediction errors, with matched likelihood conditions.}
Core-conditioned CDF distances span 0.0067--0.0450, while full-support
distances span 0.0229--0.4611. The common-shape limit ratios span
0.443--1.191 even when location and width are supplied. Thus a single
full standardized shape is not supported as a drop-in replacement.
A shared core with mass-dependent core fraction and tails is better
motivated by these diagnostics. For inspection, 181 exploratory common-shape
templates are saved at 1~MeV spacing over 60--240~MeV using the analytic
center and width laws; they are not adopted as the nominal signal model.

\clearpage
\section{Left and right signal probability and asymmetric blinding}
Asymmetry is measured around the fitted core, retaining the full
stored-histogram normalization. The upper panel measures probability
beyond $\pm2.25\sigma_{\rm nominal}$; the lower panel compares the two
halves inside that interval. Source overflow is below 0.016\% and is
excluded by the explicitly declared 0--400~MeV normalization here.
\fig{.43}{../figures/left_right_MC_probability.pdf}{Mass dependence of
the left and right selected-MC contributions around the core.}
The left external tail is larger at 120--240~MeV, whereas the right
external tail is larger at 60--100~MeV. The direction is therefore not
universal across the accepted samples.

Five windows are declared around the same measured core before fitting
the observed data: $[-2.25,+2.25]$, $[-2.5,+2.0]$, $[-2.0,+2.5]$,
$[-2.75,+2.25]$ and $[-2.25,+2.75]$, all in nominal sigma.
The redistributed pair keeps total width 4.5 sigma; each is also exactly
a shift of the window midpoint by $\mp0.25$ sigma, so it cannot isolate
tail asymmetry from center placement. The expanded pair adds half a
sigma on either side as matched controls.

Each window defines both the fitted bins and excluded GP training bins.
The GP prediction and covariance are recomputed; archived kernel,
resolution and coupling conventions remain anchored at the generated
mass. Gaussian and direct-MC templates share the same mask and background.
Actual bin counts are saved because equal continuous widths need not
select the same number of discrete observed bins.

\clearpage
\subsection{Observed asymmetric-window comparisons}
\begin{table}[H]\centering\scriptsize
\begin{tabular}{lrrrr}
\toprule
Window in nominal sigma & MC UL / symmetric & Med. $|MC/G-1|$ (\%) & Min. mass & Min. local p0 \\
\midrule
$[-2.25,+2.25]$ & 1.000--1.000 & 17.8 & 80 & 0.00787 \\
$[-2.5,+2]$ & 0.942--1.038 & 19.4 & 80 & 0.00671 \\
$[-2,+2.5]$ & 0.964--1.119 & 18.6 & 80 & 0.00854 \\
$[-2.75,+2.25]$ & 0.983--1.129 & 20.5 & 80 & 0.00761 \\
$[-2.25,+2.75]$ & 0.997--1.133 & 18.5 & 80 & 0.00868 \\
\bottomrule
\end{tabular}

\caption{Prespecified native-mass window comparisons. The minimum local
probability is descriptive and is not corrected for testing several
windows. No prescription is chosen from the observed result.}
\end{table}
\fig{.48}{../figures/asymmetric_comparisons.pdf}{Native MC limit ratios,
signed-root changes, local probabilities and fitted signal fractions for
all five windows. Full MC/Gaussian limits and probabilities are retained
in the accompanying 100-row fit ledger. Changes include altered data bins
and altered background interpolation, as well as signal acceptance.}
These are conditional observed diagnostics. A preferred production
window would require signal-recovery and coverage checks, including
signal leakage into the training sidebands, rather than selection by
the largest observed local significance or tightest upper limit.
